\section{先选工作方式，再选工具}
\label{sec:ai-entry}

\textbf{普通聊天}适合问概念、比较思路、讨论需求，或针对一小段代码请它逐步讲解；这时你主要阅读回答，再决定下一步。需要查资料、读取实验记录、分析数据并制作报告时，可以选择支持\textbf{多步骤工具任务}的工作方式，让 AI 调用工具、读取结果后继续处理。需要进入源码工程，搜索文件、修改代码、执行命令并读取报错时，则适合具备相应工具的\textbf{代码代理（Coding Agent）}。代码代理的价值来自“查看工程—采取行动—验证结果”的循环，而不只是一次生成更长的回答。

这些工作方式没有绝对互斥的边界。同一产品可以在聊天中启用搜索、文件或代码工具，代码代理也能只讨论而不修改。判断时应查看\textbf{当前启用的工具、执行环境和权限}，不能只看产品叫“聊天”还是“代理”。在本例软件中，ChatGPT 的 \textbf{Chat} 用于对话，\textbf{Work} 用于多步骤资料与产物任务，\textbf{Codex} 提供工程开发入口；这些菜单名称并不是所有产品通用的分类。

\textbf{命令行客户端（Command Line Interface，CLI）描述的是使用入口，不是另一种模型能力。}例如 Codex CLI 是使用 Codex 的终端入口。桌面客户端适合同时看对话、文件预览、浏览器和代码差异；CLI 适合已经在终端、SSH 服务器或开发容器中工作的你。能否执行代码任务，仍取决于连接的代理、主机工具、环境、权限和资料。

\begin{center}
\begin{tabular}{p{0.23\textwidth}p{0.32\textwidth}p{0.35\textwidth}}
\toprule
通用角色与本例入口 & 适合眼前什么事 & 你应当检查什么 \\
\midrule
普通聊天，例如 ChatGPT / Chat & 解释一行代码、比较两种写法、读懂报错 & 推理和示例是否正确，能否自己复述 \\
多步骤工具任务，例如 ChatGPT / Work & 查资料、分析实验记录、制作报告和图表 & 来源、数据口径、产物内容与版面 \\
代码代理，例如桌面版中的 Codex & 在工程中修复问题，同时查看差异与运行界面 & 修改范围、代码差异、测试和实际行为 \\
CLI 客户端，例如 Codex CLI & 在终端、SSH 或容器内使用代码代理 & 当前目录、运行主机、环境、命令和差异 \\
\bottomrule
\end{tabular}
\end{center}

\begin{example}[同一个问题，按任务选择入口]
假设你按下运行，计分器仍显示旧结果。先在手边的聊天工具里提供期望、实际输出和相关代码，请它解释“保存、构建、启动”分别会做什么。你可以一边读回答，一边在 IDE 里核对构建日志和启动目标。

如果要让 AI 自己查看练习副本里的文件和构建记录，就需要一个有相应文件和命令权限的代码代理。先请它只读项目，指出实际看过的文件和记录。可以从桌面操作，也可以在熟悉终端后选择 CLI。若任务只是汇总三次运行记录，用支持读取材料和生成图表的工具即可，最后要核对图表和原始数字。
\end{example}

\begin{exercise}\label{ex:ai-desktop-entry}
分别为“解释计分函数的一行代码”“在练习副本中定位一次构建错误”“把三次实验的日志整理成图表”选择工作方式与合适的客户端，可以使用本例软件或已有的其他工具，并写出各自最重要的一条验收标准。说明为什么 CLI 不是另一种模型能力，为什么不能只凭产品名称判断能否操作工程。
\end{exercise}
